\documentclass[11pt]{amsart}
\usepackage[T1]{fontenc}
\usepackage{amsmath,amssymb,amsthm}
\usepackage[margin=1.1in]{geometry}
\usepackage{booktabs}
\usepackage{enumitem}
\usepackage{hyperref}

\newtheorem{theorem}{Theorem}[section]
\newtheorem{lemma}[theorem]{Lemma}
\newtheorem{proposition}[theorem]{Proposition}
\newtheorem{corollary}[theorem]{Corollary}
\newtheorem{conjecture}[theorem]{Conjecture}
\theoremstyle{definition}
\newtheorem{definition}[theorem]{Definition}
\newtheorem{remark}[theorem]{Remark}
\newtheorem{example}[theorem]{Example}

\newcommand{\pr}{p_r}
\newcommand{\CB}{\mathrm{CB}}
\newcommand{\lean}[1]{\texttt{\hyphenchar\font=`\_ #1}}
\usepackage{array}

\title{The two-end test on cones: rooted Gallai for Eulerian graphs}
\author{Gallai project}
\date{27 November 2026}

\begin{document}
\sloppy

\begin{abstract}
The spine reduction of the project turns Gallai's path decomposition conjecture into two deletion
tests on $3$-connected graphs. We study the first test, (T1), on cones $J=B+z$. Then (T1) is exactly
the following statement RE: every $2$-connected Eulerian graph $B$ of even order $n$ has, for every
vertex $r$, a decomposition into at most $n/2$ paths of which at least two end at $r$. The spine needs
RE for such $B$ without a matching cut: at the apex of a cone over a graph of even order, the
instances of the spine are exactly RE. Gallai's conjecture implies RE, and RE implies Gallai's
conjecture for $2$-connected Eulerian graphs of even order. The correspondence is machine-checked.

Floor-bound theorems in the literature give RE for planar graphs, graphs of treewidth at most $3$,
$2$-degenerate graphs and $K_5$-minor-free graphs; the last rests on an unrefereed 2026 preprint. A doubling argument added here derives RE for
maximum degree at most $4$ from Bonamy and Perrett's ceiling theorem. We prove a lower bound
$\CB(B,r)$ on the rooted path number by counting ends. We also prove reductions at vertices of
degree $2$: in a smallest counterexample to the parity-free form RE$^*$, every degree-$2$ vertex other
than $r$ is a triangle ear, no two are twins, and at odd order there are none.

By exhaustive computation, RE holds for all $2$-connected Eulerian graphs of even order at most $10$
($301{,}882$ roots). It also holds on several families of order $12$ to $16$, including all
$258{,}774$ dense $2$-connected Eulerian graphs of order $12$ with at least $45$ edges. At order at
most $10$ the rooted path number equals the counting bound at $99.9\%$ of the roots, and every tight
root is tight for counting reasons. This leads to a strengthened conjecture UC: unless the counting
bound forbids it, one path is always to spare.
\end{abstract}

\maketitle

\section{Introduction}

All graphs are finite and simple. A \emph{decomposition} of a graph $G$ is a partition of $E(G)$ into
paths of positive length. For a decomposition $D$ and a vertex $v$, $e_D(v)$ is the number of paths of
$D$ with an end at $v$. A path has two distinct ends, and it contains $v$ at most once. Every path
through $v$ uses two edges at $v$, and every path ending at $v$ uses one, so
\begin{equation}\label{eq:par}
d(v) = e_D(v) + 2\,\#\{P\in D: v \text{ interior to } P\}, \qquad e_D(v)\equiv d(v) \pmod 2 .
\end{equation}
The path number $p(G)$ is the least size of a decomposition. For a vertex $r$, the \emph{rooted path
number} $\pr(G)$ is the least size of a decomposition with $e_D(r)\ge 1$. A graph is \emph{Eulerian}
if every degree is even; connectivity is stated separately. For Eulerian $G$, \eqref{eq:par} gives
$e_D(r)\ge 1 \iff e_D(r)\ge 2$.

\begin{definition}
\textbf{RE}$(B,r)$: $\pr(B)\le |B|/2$, for a $2$-connected Eulerian graph $B$ of even order and $r\in
V(B)$. \textbf{RE$^*$}$(B,r)$: $\pr(B)\le\lceil |B|/2\rceil$, for a connected Eulerian graph $B$ with
$|B|\ge 3$ and $r\in V(B)$.
\end{definition}

The spine of the project (Lean \lean{gallai\_iff\_spineTests'}) states that Gallai's conjecture is
equivalent to two tests on $3$-connected graphs $J$ at a vertex $z$, in their weakest form only when
$J-z$ has no matching cut. The first test, (T1) (Lean \lean{Test1}), asks: if $z$ is the only even
vertex of $J$, then for every $r\in N(z)$ the graph $J-z$ has a decomposition into at most
$(|J|-1)/2$ paths with at least two ends at $r$. Synthesis \S26 observed that on cones $J=B+z$ this is
RE. The triangle, forest and forest-through-$r$ lifts of \texttt{direct-test1} all fail on the cone
over $K_{2,6}$: there is no triangle data, the forest base fails (\texttt{direct-test1}, Example~4.7),
and so does the variant through $r$ (synthesis \S25, addendum). So the cone case needs its own
treatment.

\subsection*{What is proved, and how}
\begin{center}\small
\begin{tabular}{@{}>{\raggedright}p{0.5\textwidth}>{\raggedright}p{0.2\textwidth}p{0.22\textwidth}@{}}
\toprule
Statement & Status & Where \\
\midrule
(T1) on cones $\iff$ RE; spine scope (Cor.~\ref{cor:spine}) & proved & Lean \lean{EulerianCone.lean} \\
Gallai $\Rightarrow$ RE$^*$ at even order; RE $\iff$ leaf test; double & proved & Lean, \S\ref{sec:lit} \\
RE on planar, tw $\le 3$, $2$-degenerate, $\Delta\le4$; $K_5$-minor-free; $B-r$ a forest & from the literature; (d) unrefereed preprint, (g) via restatements & Thm.~\ref{thm:lit} \\
Counting bound $\pr\ge\CB$ & proved & Prop.~\ref{prop:cb} \\
Reductions at degree $2$; smallest counterexample, parts (a)--(c) & proved, data-checked; insertion lemma in Lean & \S\ref{sec:red}, \lean{CycleInsert.lean} \\
Smallest counterexample has $n\ge11$ (Thm.~\ref{thm:min}(d)) & computation & \S\ref{sec:comp} \\
RE, $|B|\le 10$, all roots & computation (exhaustive) & \S\ref{sec:comp} \\
RE, families of order $12$--$16$ & computation (scope in Table~\ref{tab:fam}) & \S\ref{sec:comp} \\
UC (Conj.~\ref{conj:uc}) & conjecture; exhaustive to order $10$ & \S\ref{sec:uc} \\
\bottomrule
\end{tabular}
\end{center}

\section{Cones and the spine}\label{sec:cones}

For a graph $B$ on $W$, the cone $J=B+z$ has vertex set $W\cup\{z\}$; $z$ is joined to every vertex
of $W$, and $B$ is induced on $W$. (Lean: \lean{fullCone B} on \lean{Option W}, apex \lean{none}; it is
the graph \lean{cone B univ} of \lean{Marked.lean}.) Write $n=|B|$.

\begin{proposition}[Lean: \lean{fullCone\_degree\_none}, \lean{fullCone\_degree\_some},
\lean{fullCone\_onlyEven\_iff}]\label{prop:par}
$d_J(z)=n$ and $d_J(a)=d_B(a)+1$ for $a\in W$. Hence $z$ is the only vertex of even degree of $J$ if
and only if $n$ is even and $B$ is Eulerian.
\end{proposition}

\begin{proof}
The neighbourhood of $z$ is $W$. The neighbourhood of $a$ is $N_B(a)\cup\{z\}$. So $d_J(a)$ is odd iff
$d_B(a)$ is even.
\end{proof}

\begin{proposition}[Lean: \lean{kConnected\_three\_fullCone\_iff}]\label{prop:conn}
$J$ is $3$-connected if and only if $B$ is $2$-connected. (Here $k$-connected includes having at least
$k+1$ vertices.)
\end{proposition}

\begin{proof}
Suppose $B$ is $2$-connected and $X\subseteq V(J)$ has $|X|\le 2$. If $z\in X$, then $J-X=B-(X-z)$ is
connected, since $|X-z|\le1$. If $z\notin X$, then $z$ survives, is adjacent to all of $J-X$, and $J-X$
is connected. Also $|J|=n+1\ge 4$.

Conversely, suppose $J$ is $3$-connected and $X\subseteq W$ has $|X|\le1$. Then $J-(X\cup\{z\})$ is
connected and equals $B-X$; a walk avoiding $z$ uses only edges of $B$. Also $n\ge 3$.
\end{proof}

\begin{proposition}[Lean: \lean{fullConeDelIso}, \lean{test1At\_fullCone\_iff},
\lean{test1\_fullCone\_iff}]\label{prop:t1}
$J-z\cong B$. For every $r\in W$, the \emph{conclusion} of (T1) at $(J,z,r)$ (Lean \lean{Test1At}:
$J-z$ has a decomposition into at most $(|J|-1)/2$ paths with $e(r)\ge2$) holds if and only if $B$
has a decomposition into at most $n/2$ paths with $e(r)\ge2$. Consequently (T1) holds at $(J,z)$ if and only if the following
holds: if $n$ is even and $B$ is Eulerian, then $B$ has such a decomposition for every $r$. When $B$
is $2$-connected this is RE$(B,\cdot)$; the equivalence itself does not use connectivity.
\end{proposition}

\begin{proof}
$(|J|-1)/2=n/2$, and decompositions are transported along the isomorphism with their sizes and
ends (Lean \lean{PathDecomp.exists\_of\_iso}, built on \lean{PathDecomp.mapInj}). The hypothesis
of (T1), that $z$ be the only even vertex, is Proposition~\ref{prop:par}.
\end{proof}

\begin{proposition}[Lean: \lean{test2\_fullCone\_apex}, \lean{test1\_fullCone\_some},
\lean{test2\_fullCone\_some\_of\_even}]\label{prop:vac}
The other tests at a cone are vacuous: (T2) at $z$ always; (T1) at every $a\in W$ always; (T2) at every
$a\in W$ when $n$ is even.
\end{proposition}

\begin{proof}
No graph has exactly one vertex of odd degree (handshake lemma).
\begin{itemize}[leftmargin=*]
\item (T2) at $(z,a)$ needs $z,a$ to be the only even vertices of $J$. Then $a$ would be the only odd
vertex of $B$.
\item (T1) at $a$ needs $a$ to be the only even vertex. Then $z$ is odd, so $n$ is odd and $|J|$ is even,
and $J$ would have the odd number $|J|-1$ of odd vertices.
\item (T2) at $(a,x)$ with $n$ even: $z$ is even, so $x=z$, and $a$ would be the only odd vertex of $B$.
\end{itemize}
\end{proof}

When $n$ is odd, (T2) at a non-apex vertex $a$ of the cone is a genuine test. It concerns
$B$ with exactly two odd vertices, which are adjacent. It is not part of RE and is not treated here.

\begin{proposition}[Lean: \lean{hasMatchingCut\_delVert\_fullCone\_iff}]\label{prop:mc}
$J-z$ has a matching cut if and only if $B$ has one.
\end{proposition}

\begin{proof}
Matching cuts are invariant under isomorphism (Lean \lean{HasMatchingCut.of\_iso}).
\end{proof}

\begin{corollary}[Lean: \lean{rootedDecomp\_of\_spineTestsNMC}, \lean{fullConeTests\_iff},
\lean{rootedDecomp\_of\_gallai}]\label{cor:spine}
\leavevmode
\begin{enumerate}[label=(\roman*)]
\item At the apex, the instances of the spine's weakest form \lean{SpineTestsNMC} on cones over
graphs of even order are exactly RE$(B,r)$, at every $r$, for $2$-connected Eulerian $B$
\emph{without a matching cut}. Lean proves the direction \lean{SpineTestsNMC} $\Rightarrow$ RE for
these $B$ (\lean{rootedDecomp\_of\_spineTestsNMC}). It also proves the unrestricted equivalence: for $B$
of even order, every test at every vertex of $B+z$ holds if and only if, when $B$ is Eulerian, $B$
has a rooted decomposition with at most $n/2$ paths at every $r$ (\lean{fullConeTests\_iff}). For
odd $n$, the spine's cone instances also include the genuine (T2) tests at non-apex vertices
(Proposition~\ref{prop:vac}), which are not RE.
\item Gallai's conjecture implies $\pr(B)\le |B|/2$ for every connected Eulerian $B$ of even order and
every $r$ with $d(r)>0$. So
\[
\text{Gallai}\ \Longrightarrow\ \text{RE}\ \Longrightarrow\ \text{Gallai for $2$-connected
Eulerian graphs of even order.}
\]
\end{enumerate}
\end{corollary}

\begin{proof}
(i) Propositions~\ref{prop:par}--\ref{prop:mc}. (ii) Gallai's conjecture is equivalent to the rooted
private-leaf test (Lean \lean{gallai\_iff\_rootedTest}): $p(B+\text{leaf at } r)\le|B|/2$ for
connected $B$ of even order. Lemma~\ref{lem:leaf} below converts this to $\pr(B)\le |B|/2$. The last
implication holds since $p(B)\le\pr(B)$.
\end{proof}

Matching cuts are common among small $2$-connected Eulerian graphs, but most have none. Among the $1$,
$7$, $162$, $30{,}054$ graphs of orders $4,6,8,10$, those without a matching cut number $0$, $5$,
$130$, $27{,}933$ (\texttt{scripts/mcut.c}). Cycles of length at least $4$ have matching cuts. $K_{2,6}$
has none.

\section{The leaf graph, the double, and the literature}\label{sec:lit}

Let $B^r$ be $B$ with a new vertex $\ell$ joined to $r$ only. Let $B\oplus_r B$ be two disjoint copies
of $B$ joined by an edge between the two copies of $r$ (Lean \lean{dbl}).

\begin{lemma}[threshold form in Lean: \lean{rootedDecomp\_iff\_addLeaf}]\label{lem:leaf}
If $d_B(r)>0$, then $\pr(B)=p(B^r)$.
\end{lemma}
In Lean, the threshold form for even $d_B(r)>0$ is stated: $B$ has at most $|B|/2$ paths with two ends at
$r$ if and only if $p(B^r)\le|B|/2$. Evenness is needed only to turn one end into two. The
equality is assembled on paper from the same pieces (\lean{pathNumber\_addLeaf\_le},
\lean{exists\_leafCost\_le}).

\begin{proof}
($\ge$) Extend a path ending at $r$ by $r\ell$.
($\le$) Let $D$ decompose $B^r$, and let $P$ be the path containing $r\ell$.
\begin{itemize}[leftmargin=*]
\item If $P=r\ell$, then $D-P$ decomposes $B$ into $|D|-1$ paths. If $r$ is an end of none of them, some
path passes through $r$: some edge at $r$ lies on a path, and $r$ is not its end. Splitting it there gives $|D|$
paths with two ends at $r$.
\item Otherwise $P-\ell$ ends at $r$, and the count is unchanged.
\end{itemize}
\end{proof}

\begin{lemma}[Lean: \lean{two\_mul\_pathNumber\_addLeaf\_le}, \lean{rootedDecomp\_of\_dbl}]\label{lem:dbl}
$2p(B^r)\le p(B\oplus_r B)+1$. Hence, if $|B|$ is even, $d_B(r)$ is even and positive, and
$p(B\oplus_r B)\le |B|$, then $\pr(B)\le |B|/2$.
\end{lemma}

\begin{proof}
Exactly one path $P$ of a decomposition of the double uses the bridge. Cutting $P$ at the bridge, and
keeping each half on its own side together with the bridge edge as the leaf edge, gives decompositions
of the two leaf graphs. Their sizes add up to at most $|D|+1$. Then
$\pr(B)=p(B^r)\le(|B|+1)/2$, which is $\le|B|/2$ since $|B|$ is even.
\end{proof}

The leaf lemma converts a floor bound $p\le\lfloor n'/2\rfloor$ on a class closed under adding a leaf
into RE. The double needs only a ceiling bound, since $|B\oplus_rB|$ is even, but the class must
contain the double.

\begin{theorem}[RE from the literature]\label{thm:lit}
Let $B$ be a connected Eulerian graph of even order and $r\in V(B)$. Then $\pr(B)\le|B|/2$ in each of
the following cases.
\begin{enumerate}[label=(\alph*)]
\item $B$ is planar.
\item $B$ has treewidth at most $3$.
\item $B$ is $2$-degenerate.
\item $B$ is $K_5$-minor-free (conditional on \cite{ZLH}, an unrefereed preprint).
\item $\Delta(B)\le 4$.
\item $\Delta(B-r)\le3$.
\item $B-r$ is a forest (via the restatements of Pyber's theorem in \cite{BS,BP}).
\end{enumerate}
\end{theorem}

\begin{proof}
In (a)--(d) and (f)--(g) we apply Lemma~\ref{lem:leaf}. The graph $B^r$ has odd order $|B|+1$ and a leaf,
so it is none of the exceptions $K_3$, $K_5^-$, $K_5$. Nor is it a SET graph (a graph whose
even-vertex subgraph is a triangle and whose every odd vertex has at least two even neighbours): the odd
vertex $\ell$ has no even neighbour. It lies in the class whenever $B$ does.
\begin{enumerate}[label=(\alph*)]
\item Blanch\'e, Bonamy and Bonichon \cite[Thm.~1.2]{BBB}: every connected planar graph except $K_3$ and
$K_5^-$ decomposes into $\lfloor n/2\rfloor$ paths.
\item Botler, Sambinelli, Coelho and Lee \cite[Thm.~3.2]{BSCL}: connected partial $3$-trees, except
$K_3$ and $K_5^-$.
\item Anto and Basavaraju \cite[Thm.~1]{AB}: connected $2$-degenerate graphs, except $K_3$.
\item Zhang, Liu and Hong \cite[Thm.~1.1]{ZLH}: connected $K_5$-minor-free graphs, except $K_3$ and
$K_5-e$. This is an arXiv preprint (November 2026), not yet refereed.
\item This case uses Lemma~\ref{lem:dbl} instead. The double is connected and has maximum degree
$\max(\Delta(B),d(r)+1)\le 5$. Bonamy and Perrett \cite[Thm.~1.3]{BP} decompose every connected graph
of maximum degree at most $5$ into $\lceil n/2\rceil$ paths. So $p(B\oplus_rB)\le|B|$.
\item Botler and Sambinelli \cite[Thm.~13]{BS}: if $G$ is connected and simple and
$\Delta(\mathrm{EV}(G))\le3$, then $G$ is a Gallai graph or a SET graph. Here $\mathrm{EV}(B^r)=B-r$.
\item Pyber's theorem: if the even vertices induce a forest, then $p\le\lfloor n/2\rfloor$. We used the
restatements in \cite[Thm.~2(b)]{BS} and \cite{BP}, not Pyber's paper itself, which we could not access
(\texttt{research/notes/eulerian-gallai-literature.md}).
\end{enumerate}
\end{proof}

The cone over $K_{2,6}$, which defeats the terminal-structure lifts, is covered by (a) and by (g).
Every class of Theorem~\ref{thm:lit} is sparse or of bounded degree. The tight instances of RE in
\S\ref{sec:comp} are dense.

No paper we found proves a rooted statement for Eulerian graphs of RE's strength. The closest are
the following:
\begin{itemize}[leftmargin=*]
\item Fan--Hou--Zhou \cite[Thm.~5]{FHZ} and Xie \cite[Thm.~1.4]{Xie} prove bounds $(n+1)/2$ with
two ends at a prescribed vertex. We read Fan--Hou--Zhou only in its abstract and in Xie's verbatim
restatement.
For $2$-connected Eulerian graphs their hypotheses force a cycle, or an edge with all other vertices of
degree $2$.
\item Botler--Sambinelli \cite[Lemma~10]{BS} is a rooted lemma for ESET graphs, a non-Eulerian class
that contains the SET graphs.
\end{itemize}
Chu and Wang \cite{CW}, whose title is \emph{Path decompositions of Eulerian graphs}, prove $3n/5$ for
Eulerian graphs whose triangles are pairwise at distance at least $3$, and for graphs with a triangle
removal set. The survey with access levels is \texttt{research/notes/eulerian-gallai-literature.md}.

\subsection*{Results in the project's own collection}
The project's collection (\texttt{gallai/}, unrefereed) contains Eulerian results of two kinds.
\begin{itemize}[leftmargin=*]
\item \emph{The dense regime, unconditionally.} Marked-reduction Cor.~168: in a connected Eulerian
graph of order $2k$, any $k$ vertices of degree $2k-2$ are the ends, twice each, of a $k$-path
decomposition. Thm.~170: if at least $k-1$ vertices have degree $2k-2$, then RE holds at
\emph{every} root. Thm.~172 extends this to $(1-1/1000)k$ such vertices for large $k$.
\begin{itemize}
\item The proof of Thm.~170 uses $6{,}131$ finite certificates whose supplement is not in the
repository. We checked both statements exactly, with the profile oracle \texttt{tools/route\_ext.c},
on every connected Eulerian graph of even order at most $10$: all $1{,}579$ instances of Cor.~168 and
all $14{,}096$ instances of Thm.~170 are realisable (\texttt{scripts/dense\_check.py}).
\item The hypothesis of Thm.~170 holds in only $319$ of these $31{,}219$ graphs. At order $10$ it
covers $1{,}323$ of the $14{,}156$ tight roots: most roots with $d(r)=n-2$ lie in graphs with fewer
than $n/2-1$ such vertices.
\end{itemize}
\item \emph{Reductions inside Gallai's induction.} Marked-reduction Thm.~177, Cor.~178 and Thm.~180,
and separating-pairs Cor.~13.12, describe a least failure of the rooted private-leaf test, which is
equivalent to Gallai's conjecture, when the failing host is Eulerian. Examples: a degree-$(n-2)$ root
whose non-neighbour has degree $2$ forces neighbours of degree at least $16$; a degree-$2$ vertex that
is not an ear forces all other degrees to be at least $18$; and, conditionally, minimum degree at
least $10$.
\begin{itemize}
\item These results assume the rooted bound for \emph{all} connected graphs of smaller even order.
They constrain a first counterexample to Gallai and do not prove RE in any range.
\item Theorem~\ref{thm:min} works in a different induction, inside the Eulerian class but over both
parities. There RE$^*$ at odd order supplies the slack, so no degree hypotheses are needed. The two
are not comparable: RE$^*$ at odd order is implied by Gallai but is not a consequence of RE.
\end{itemize}
\end{itemize}


\section{A lower bound by counting ends}\label{sec:cb}

\begin{proposition}\label{prop:cb}
Let $B$ be Eulerian and $r\in V(B)$, and let $D$ be a decomposition into $p$ paths with $e_D(r)\ge1$.
Then the numbers $e_v=e_D(v)$ are even and satisfy
\[
0\le e_v\le c_v(p):=\min(d_v,\,2p-d_v),\qquad e_r\ge 2,\qquad \textstyle\sum_v e_v=2p .
\]
Hence $\pr(B)\ge\CB(B,r)$, the least $p$ such that $c_r(p)\ge2$, $c_v(p)\ge 0$ for all $v$, and
$\sum_v c_v(p)\ge 2p$.
\end{proposition}

\begin{proof}
By \eqref{eq:par}, $v$ lies on exactly $(d_v-e_v)/2+e_v=(d_v+e_v)/2$ paths, which is at most $p$. So
$e_v\le 2p-d_v$. Clearly $e_v\le d_v$. All quantities are even because every $d_v$ is even. Conversely,
if the three conditions hold, even values $e_v$ with these bounds and sum $2p$ exist. So $\CB$ is
exactly the least $p$ for which these necessary conditions can be met.
\end{proof}

Three special cases give the three classical lower bounds:
\begin{itemize}[leftmargin=*]
\item $c_r\ge2$ gives $p\ge d(r)/2+1$;
\item $c_v\ge0$ gives $p\ge\Delta/2$;
\item $\sum_v(2p-d_v)\ge2p$ gives $p\ge\lceil m/(n-1)\rceil$.
\end{itemize}
The bound $\CB$ can exceed all three.

\begin{example}\label{ex:ds}
Let $B$ be $K_8$ minus a double star: the centres $s,t$ are adjacent in the complement, with four and
two leaves (graph6 \texttt{GTm\~{}vs}). Then $d(s)=2$, $d(t)=4$, all other degrees are $6$, and $m=21$.
At $p=3$ we get $\sum_v c_v=2+2+6\cdot0=4<6$, so $\CB=4=|B|/2$ at every root, while the three classical
bounds give $3$ at the roots $s$ and $t$. Directly: three paths with $21=3\cdot 7$ edges are Hamiltonian,
and all three would visit $s$, which lies on at most $d(s)=2$ paths.
\end{example}

\section{Reductions at vertices of degree two}\label{sec:red}

Throughout, $B$ is connected and Eulerian with $n=|B|\ge 4$, and $r\in V(B)$. A vertex $v$ of degree $2$
with $N(v)=\{a,b\}$ is an \emph{ear} if $ab\in E(B)$.

\begin{lemma}[suppression]\label{lem:supp}
Let $v\ne r$ have degree $2$ with $N(v)=\{a,b\}$ and $ab\notin E(B)$. Then $B/v:=B-v+ab$ is a
connected Eulerian simple graph on $n-1$ vertices, and $\pr(B)\le\pr(B/v)$.
\end{lemma}

\begin{proof}
Degrees are unchanged, and connectivity is clear. In a decomposition of $B/v$, replace the edge $ab$
by $avb$ inside its path. The vertex $v$ is new, so this is still a path, and no end changes.
\end{proof}

\begin{lemma}[cycle insertion; Lean: \lean{PathDecomp.cycle\_insert}]\label{lem:ins}
Let $G'$ be a graph, and let $a\in V(G')$ with $d_{G'}(a)>0$. Let $C=a\,c_1\cdots c_l\,a$ ($l\ge2$) be a
cycle that is edge-disjoint from $G'$, such that $c_l\notin V(G')$ and at most one of
$c_1,\dots,c_{l-1}$ lies in $V(G')$. For every decomposition $D'$ of $G'$, the graph $G=G'\cup C$ has
a decomposition $D$ with:
\begin{itemize}
\item $|D|=|D'|+1$;
\item $e_D(x)=e_{D'}(x)$ for every $x\in V(G')$;
\item $e_D(c_l)=2$, and $e_D=0$ at the other new vertices.
\end{itemize}
\end{lemma}

\begin{proof}
\emph{Case 1: $a$ is interior to some $Q\in D'$.} At most one vertex $c_j$ of $C$ other than $a$ lies on
$Q$, and it lies on one side of $a$. Orient $Q=Q_1\,a\,Q_2$ with $c_j\notin Q_1$; both $Q_1$ and $Q_2$
are nonempty. Replace $Q$ by $Q_1\,a\,c_1\cdots c_l$ and $c_l\,a\,Q_2$. The first is a path because the
$c_i$ avoid $Q_1$. The second is a path because $c_l$ is new.

\emph{Case 2: otherwise.} Then $a$ is an end of some path $P=a\,P'$. Replace $P$ by $c_l\,a\,P'$ and add
the path $c_l\,c_{l-1}\cdots c_1\,a$.

In both cases every edge of $C$ is used once. The old ends stay where they were, except that in Case~2
the end of $P$ at $a$ moves to $c_l$ and the new path ends at $a$. So $c_l$ gains two ends.

In the Lean statement, $C$ is given as a path $R\colon a\to c_l$ of $G$ together with the edge $c_la$, and
``$x\notin V(G')$'' is weakened to ``$x$ has no edge in $G'$''. Case~2 is used whenever $a$ is an end of
some path, and Case~1 otherwise.
\end{proof}

\begin{proposition}[ears; the insertion step in Lean: \lean{PathDecomp.ear\_insert}]\label{prop:ear}
Let $v\ne r$ be an ear with $N(v)=\{a,b\}$, and suppose $B':=B-v-ab$ is connected. Then
$\pr(B)\le\pr(B')+1$. In particular RE$^*(B',r)$ implies RE$^*(B,r)$ when $n$ is odd.
\end{proposition}

\begin{proof}
$B'$ is connected, Eulerian, and has $n-1\ge3$ vertices, so $d_{B'}(a)>0$. Apply Lemma~\ref{lem:ins} with
$C=a\,b\,v\,a$ ($c_1=b$, $c_2=v$) to an optimal rooted decomposition of $B'$. For odd $n$,
$\lceil (n-1)/2\rceil+1=\lceil n/2\rceil$.
\end{proof}

\begin{proposition}[ears at a cut]\label{prop:earcut}
Let $v\ne r$ be an ear with $N(v)=\{a,b\}$, and suppose $B-v-ab$ is disconnected. Then $B-v$ is
connected, and $B-v-ab$ has exactly two components $B_a\ni a$ and $B_b\ni b$. Each is Eulerian, and is
either a single vertex or has at least $3$ vertices. Choose the labels so that $r\in V(B_a)$.
\begin{enumerate}[label=(\roman*)]
\item If $|B_a|,|B_b|\ge3$, then $\pr(B)\le p_r(B_a)+p_b(B_b)$.
\item If $B_b=\{b\}$, then $\pr(B)\le p_r(B_a)+1$.
\item If $B_a=\{a\}$, so that $r=a$, then $\pr(B)\le p_b(B_b)+1$.
\end{enumerate}
In every case RE$^*$ for the smaller graphs gives RE$^*(B,r)$.
\end{proposition}

\begin{proof}
$B-v$ is connected because $a$ and $b$ are adjacent. Deleting the edge $ab$ leaves at most two
components. An Eulerian connected graph on two vertices does not exist.

(i) Take optimal decompositions $D_a$ (rooted at $r$) and $D_b$ (rooted at $b$), and a path $P_b=P_b'\,b\in
D_b$ ending at $b$.
\begin{itemize}[leftmargin=*]
\item If $a$ is interior to some $Q=Q_1\,a\,Q_2\in D_a$, replace $Q,P_b$ by $Q_1\,a\,v\,b$ and
$P_b'\,b\,a\,Q_2$.
\item Otherwise $a$ is an end of some $P_a=a\,P_a'\in D_a$. Replace $P_a,P_b$ by $P_b'\,b\,a\,P_a'$ and
add $a\,v\,b$.
\end{itemize}
These are paths since $V(B_a)\cap V(B_b)=\emptyset$ and $v$ is new. The number of paths is
$|D_a|+|D_b|$, and every old vertex keeps its number of ends.

(ii) Here $b$ has neighbours $a,v$ only. Apply Lemma~\ref{lem:ins} to $D_a$ with $C=a\,b\,v\,a$.

(iii) Apply Lemma~\ref{lem:ins} to $D_b$ with $C=b\,v\,a\,b$. The new vertex $a=r$ is $c_l$, so it gets
two ends.

Bounds: $n=|B_a|+|B_b|+1$. In (i), $\lceil|B_a|/2\rceil+\lceil|B_b|/2\rceil\le\lceil n/2\rceil$ by
parity. In (ii) and (iii), $\lceil (n-2)/2\rceil+1=\lceil n/2\rceil$.
\end{proof}

\begin{proposition}[twins]\label{prop:twin}
Let $u\ne v$, both different from $r$, have degree $2$ and $N(u)=N(v)=\{a,b\}$. Then RE$^*$ for smaller
graphs gives RE$^*(B,r)$.
\end{proposition}

\begin{proof}
If $ab\notin E$, use Lemma~\ref{lem:supp}. If $ab\in E$, then $B''=B-u-v$ is connected (it contains $ab$),
Eulerian, and has at least $3$ vertices: $d(a)\ge 4$, since $a$ has neighbours $u,v,b$ and even degree.
Lemma~\ref{lem:ins} with $C=a\,u\,b\,v\,a$ ($c_2=b$ is the one old vertex) gives
$\pr(B)\le\pr(B'')+1\le\lceil(n-2)/2\rceil+1=\lceil n/2\rceil$.
\end{proof}

\begin{proposition}[the root of degree two, odd order]\label{prop:root}
Let $n$ be odd, $d(r)=2$ and $N(r)=\{a,b\}$. If $ab\notin E$, or if $ab\in E$ and $B-r-ab$ is connected,
then RE$^*$ for smaller graphs gives RE$^*(B,r)$.
\end{proposition}

\begin{proof}
If $ab\notin E$: $B/r$ is connected Eulerian of even order $n-1$, so $p(B/r)\le\lceil(n-1)/2\rceil$.
Subdivide $ab$ by $r$ and split the path through $r$ at $r$. This gives
$\pr(B)\le (n-1)/2+1=\lceil n/2\rceil$.
If $ab\in E$: apply Lemma~\ref{lem:ins} to $B-r-ab$ at $a$ with $C=a\,b\,r\,a$. The vertex $r=c_l$ gets two
ends, so $\pr(B)\le p(B-r-ab)+1\le\lceil n/2\rceil$.
\end{proof}

\begin{theorem}[a smallest counterexample]\label{thm:min}
Let $(B,r)$ be a counterexample to RE$^*$ with $|V(B)|+|E(B)|$ minimum. Then:
\begin{enumerate}[label=(\alph*)]
\item every vertex $v\ne r$ of degree $2$ is an ear, with $N(v)=\{a,b\}$ and $B-v-ab$ connected;
\item no two vertices other than $r$ of degree $2$ have the same neighbourhood;
\item if $n$ is odd, every vertex other than $r$ has degree at least $4$. If moreover $d(r)=2$, then
$N(r)=\{a,b\}$ is an edge and $B-r-ab$ is disconnected;
\item $n\ge 11$ (by computation, \S\ref{sec:comp}).
\end{enumerate}
\end{theorem}

\begin{proof}
RE$^*$ holds for $K_3$, so $n\ge4$. Each of Lemma~\ref{lem:supp} and Propositions~\ref{prop:ear}--\ref{prop:root}
produces smaller connected Eulerian graphs, with at least $3$ vertices or trivial, that satisfy
RE$^*$ by minimality, and it gives RE$^*(B,r)$:
\begin{itemize}[leftmargin=*]
\item (a): Lemma~\ref{lem:supp} with $\lceil(n-1)/2\rceil\le\lceil n/2\rceil$, and
Proposition~\ref{prop:earcut};
\item (b): Proposition~\ref{prop:twin};
\item (c): Proposition~\ref{prop:ear} and (a), and Proposition~\ref{prop:root};
\item (d): computation (\S\ref{sec:comp}): RE$^*$ holds for every connected Eulerian graph of order
$3\le n\le10$, at every root.
\end{itemize}
\end{proof}

\begin{remark}[what is left at degree two]\label{rem:even}
The remaining case is an ear $v$ at even order with $B'=B-v-ab$ connected. Then $B'$ has odd order,
and RE$^*(B')$ allows $n/2$ paths, so Lemma~\ref{lem:ins} gives $n/2+1$, one too many. The insertion
is free in either of two situations:
\begin{itemize}[leftmargin=*]
\item $a$ is an end of two paths $P_1,P_2$ of $D'$ with $b\notin P_1$: replace them by $P_1\,a\,b\,v$
and $P_2\,a\,v$. This moves two ends from $a$ to $v$, so it keeps an end at $r$ only if $a\ne r$ or
$e_{D'}(a)\ge4$;
\item $a,b$ are ends of vertex-disjoint paths $P_a,P_b$: merge them through $ab$ and add $a\,v\,b$.
\end{itemize}
These are avoidance conditions of the kind that the project's open core (L1$^*$, DL) is about. On the
data, the insertion bound $\pr(B')+1$ equals $n/2+1$ in $56$ of the $1{,}225$ even-order ear instances
of order at most $9$, because $\pr(B')=n/2$ there. Yet $\pr(B)\le n/2$ always held.
\end{remark}

\begin{remark}[fat thetas]\label{rem:theta}
The configuration that defeats the terminal-structure lifts of \texttt{direct-test1} consists of many
degree-$2$ vertices of $B$ between the root and a second hub, as in $K_{2,t}$. The three order-$9$
roots missed by all lifts have this shape (\texttt{direct-test1}, Example~4.7, and the triangle-base
data). Two such vertices are twins with $N=\{r,r'\}$. By Lemma~\ref{lem:supp}, or by
Proposition~\ref{prop:twin} when $rr'\in E$, they cannot both be present in a smallest counterexample
to RE$^*$: Theorem~\ref{thm:min}(b). So the fat theta is reducible in the RE$^*$ induction,
whichever hub is the root.
\end{remark}


\emph{Data check.} \texttt{scripts/reductions\_check.py} executes every construction of this section
on optimal decompositions produced by the solver, and verifies each output independently. It checks
the partition into paths, the size, and the end at $r$ in every case. It checks the ends at every old
vertex for the ear and twin insertions, the only claims that need them. It checks the bound $\lceil n/2\rceil$
wherever the text claims it.
\begin{itemize}[leftmargin=*]
\item Every connected Eulerian graph of order $4$ to $9$, every root and every applicable vertex:
$15{,}776$ suppressions, $13{,}827$ ear insertions, $2{,}537$ ears at a cut, $3{,}559$ twin
insertions, $2{,}024$ root suppressions and $1{,}763$ root-ear insertions, with no failure
(\texttt{data/out/reductions\_check\_le9.txt}).
\item At order $10$: the first $3{,}000$ of the $2$-connected graphs and all $972$ connected graphs
that are not $2$-connected. That is $62{,}838$ suppressions, $54{,}306$ ear insertions,
$16{,}002$ ears at a cut, $19{,}472$ twin insertions, $6{,}982$ root suppressions and $6{,}034$ root-ear
insertions, with no failure (\texttt{reductions\_check\_10.txt}).
\end{itemize}

\section{Computation}\label{sec:comp}

\subsection{The solver and its validation}
\texttt{scripts/re.c} decides $\pr(B)\le t$ by depth-first search over paths. The first path starts at
$r$. Afterwards, a vertex of odd degree in the remaining edge set $R$ must be an end of some path, so
the search branches over the paths starting there; if there is none, it branches over the paths
through the lowest edge of $R$. Both branchings are complete. Components of $R$ are solved separately,
and failures are memoised.

The lower bound used for pruning is $\max(1,\mathrm{odd}/2,\max_v\lceil d_v/2\rceil,\lceil
|E|/(|V|-1)\rceil)$ per component. All of its terms are valid bounds for arbitrary graphs.
The validation has four parts:
\begin{itemize}[leftmargin=*]
\item \textbf{Brute force.} \texttt{scripts/bf.py} enumerates all paths and all exact covers. On all
$473$ connected graphs of order $3$ to $7$ with at most $10$ edges it agrees with \texttt{re.c} on all
$3{,}622$ values of $\pr$ (every root) and $p$.
\item \textbf{Witnesses.} Every upper bound below comes with a witness checked by
\texttt{scripts/verify\_witness.py} or, in UC mode, \texttt{scripts/verify\_uc.py}. The full witness dumps
are regenerated by \texttt{run\_all.sh}; the shipped tables keep the sizes only.
\item \textbf{Lower bounds.} Every lower bound used below is either the counting bound of
Proposition~\ref{prop:cb}, which is a proof, or was re-decided by CP-SAT. The CP-SAT model is the
project's circuit model \texttt{tools/pathnum2.py}, which shares no code with \texttt{re.c}. It agrees on
all $219$ roots with $\pr>\CB$ (\texttt{data/out/above\_cb\_cpsat.txt}). It also agrees on the $6$ UC failures
of \S\ref{sec:uc} (\texttt{ucfail\_cpsat.txt}).
\item \textbf{Generation.} The graphs come from nauty's \texttt{geng} (Debian package). The counts of
connected Eulerian graphs of orders $3,\dots,10$ are $1,1,4,8,37,184,1782,31026$.
\end{itemize}

\subsection{Exhaustive results}
\begin{table}[h]
\begin{tabular}{@{}rrrrrrr@{}}
\toprule
$n$ & graphs & roots & $\pr\le n/2$ & tight ($\pr=n/2$) & $\pr=\CB$ & $\pr>\CB$ \\
\midrule
4 & 1 & 4 & 4 & 4 & 4 & 0\\
6 & 7 & 42 & 42 & 23 & 42 & 0\\
8 & 162 & 1{,}296 & 1{,}296 & 242 & 1{,}289 & 7\\
10 & 30{,}054 & 300{,}540 & 300{,}540 & 14{,}156 & 300{,}328 & 212\\
\bottomrule
\end{tabular}
\caption{All $2$-connected Eulerian graphs of even order $n\le10$. Every tight root has
$\CB(B,r)=n/2$.}\label{tab:ex}
\end{table}

At order $10$ the solver's budget left $8$ minima uncertified, although RE was witnessed at each of
them. Seven have $\CB=5=n/2$, so they are tight. At the eighth (\texttt{ITzvn\^{}nzw}, $r=4$),
$\CB=4$; an unbudgeted run and CP-SAT both found $4$ paths, so $\pr=\CB=4$.

At order $10$ the $14{,}156$ tight roots are tight for three counting reasons:
\begin{itemize}[leftmargin=*]
\item $14{,}131$ have $d(r)=n-2$, so the bound $d(r)/2+1$ is already $n/2$;
\item $20$ others meet the edge bound $\lceil m/9\rceil=5$;
\item the remaining $5$ need the full bound $\CB$, as in Example~\ref{ex:ds}. They are the roots
$0,1,2$ of \texttt{IEl\}\}\~{}\^{}vw} and $1,7$ of \texttt{ITm||\~{}\}\~{}g}, with $m=35,36$.
\end{itemize}
At order $8$, two tight roots need the full bound: those of Example~\ref{ex:ds}.

\emph{All connected Eulerian graphs, both parities.} RE$^*$ holds at all $17{,}844$ roots of the
connected Eulerian graphs of orders $3$ to $9$. It also holds at all $9{,}720$ roots of the $972$
connected but not $2$-connected Eulerian graphs of order $10$. Together with Table~\ref{tab:ex} this
gives Theorem~\ref{thm:min}(d).

\subsection{Families of order 12 to 16}
\begin{table}[h]
\begin{tabular}{@{}lrrl@{}}
\toprule
family ($2$-connected Eulerian) & graphs & roots & result\\
\midrule
order $12$, $m\ge45$ (complement $\le21$ edges, all odd) & 258{,}774 & 3{,}105{,}288 & RE; UC \\
order $12$, $4$-, $6$-, $8$-, $10$-regular & 1{,}542; 7{,}849; 94; 1 & 113{,}832 & RE \\
order $14$, $4$-regular & 88{,}126 & 1{,}233{,}764 & RE \\
order $14$, $10$-regular & 540 & 7{,}560 & RE \\
complete multipartite, even parts, $n\le16$ & 58 & 760 & RE \\
random, orders $12$, $14$, $16$ & 2{,}100 samples each & 88{,}200 & RE \\
\bottomrule
\end{tabular}
\caption{Families. RE was decided at every root (witness verified). For the dense order-$12$ family
UC (\S\ref{sec:uc}) was also decided at every root. $32$ roots needed CP-SAT, and all have
$\pr=\CB=5$. The random samples may contain isomorphic repeats.}\label{tab:fam}
\end{table}

Families with maximum degree $4$ are covered by Theorem~\ref{thm:lit}(e). They serve here as a solver
check. The dense family contains every $2$-connected Eulerian graph of order $12$ with at least $45$
edges, so it contains every graph of that order whose edge bound $\lceil m/11\rceil$ is at least
$5=n/2-1$. The slice $m=45$ was added after audit A30, whose own run on it agrees. Tight roots also occur in sparse graphs: for instance $K_{2,10}$ is tight at its two roots of
degree $10=n-2$.

\section{Tightness and a strengthened conjecture}\label{sec:uc}

\begin{conjecture}[UC]\label{conj:uc}
For every $2$-connected Eulerian graph $B$ and every $r\in V(B)$,
\[
\pr(B)\le\max\bigl(\CB(B,r),\ \lceil |B|/2\rceil-1\bigr).
\]
\end{conjecture}

For even order, UC strengthens RE. Since $\CB\le n/2$ always, UC says that one path is to spare
unless counting forbids it. For odd order, UC is a rooted Eulerian analogue of Bonamy and
Perrett's question \cite[Question~1.1]{BP}, whether every connected graph that is not an odd
semi-clique has $p\le\lfloor n/2\rfloor$. The counting bound takes the place of the odd
semi-cliques, and it explains $K_3$ and $K_5$.

\emph{Evidence.} UC holds at every root of every $2$-connected Eulerian graph of order $3$ to $10$,
odd orders included, and at every root of the dense order-$12$ family. It needs $2$-connectivity: it
fails at $6$ roots of connected graphs with a cut vertex ($1$ at order $8$, $5$ at order $10$). An
example is \texttt{GCXnbW}: a triangle attached at a cut vertex to a $6$-vertex block, rooted at the
cut vertex, where $\pr=4=n/2>\CB=3$. RE$^*$ still holds there.

\begin{remark}[roots of degree $n-2$]\label{rem:degn2}
At order $10$, $99.8\%$ of the tight roots have $d(r)=n-2$.
\begin{itemize}[leftmargin=*]
\item Then every one of the $n/2$ paths meets $r$, and exactly two end there.
\item $B-r$ has exactly one even vertex, the non-neighbour $s$ of $r$.
\item Deleting $r$ cuts each path into one or two \emph{pieces}, possibly trivial. There is one piece for
each neighbour $x$ of $r$, the part of the path containing $rx$ on the side of $x$.
\end{itemize}
So RE at such a root is a statement about $B-r$, a graph with one even vertex. It asks for a family
$(Q_x)_{x\in N(r)}$ with the following properties:
\begin{itemize}[leftmargin=*]
\item each $Q_x$ is a path of $B-r$ with an end at $x$, or the trivial path $x$;
\item the nontrivial $Q_x$ partition $E(B-r)$;
\item $N(r)$ splits into $n/2-2$ pairs $\{x,y\}$ with $Q_x,Q_y$ vertex-disjoint, and two singletons.
\end{itemize}
Conversely, every such family reassembles into a rooted decomposition of $B$ with $n/2$ paths:
$Q_x\,x\,r\,y\,Q_y$ for each pair, and $Q_x\,x\,r$ for each singleton.
This is a constrained Lov\'asz decomposition problem with prescribed attachments, of the
terminal-matching type studied elsewhere in the project. It is the natural next target for the cone
case of (T1).
\end{remark}

\section*{Files}
\begin{itemize}[leftmargin=*]
\item Audit: \texttt{research/audit/A30-eulerian-cone.md} (adversarial referee). It found the proofs of
\S\S2--5 sound and the Lean faithful; its findings (one false coverage sentence, one gap in a remark,
wording and provenance) have been applied, as listed at the end of the report.
\item Lean: \texttt{lean/GallaiLean/EulerianCone.lean} (\S\ref{sec:cones}, Lemmas~\ref{lem:leaf}, \ref{lem:dbl}) and
\texttt{lean/GallaiLean/CycleInsert.lean} (Lemma~\ref{lem:ins}, the insertion step of Proposition~\ref{prop:ear}). No \texttt{sorry}. Every theorem named above
depends only on \texttt{propext}, \texttt{Classical.choice} and \texttt{Quot.sound}, and the full
library builds.
\item Scripts: \texttt{scripts/run\_all.sh} regenerates every number above.
\begin{itemize}
\item \texttt{re.c} is the solver; compile with \texttt{-DWIDE} for up to $128$ edges.
\item \texttt{bf.py}, \texttt{compare\_bf.py}, \texttt{verify\_witness.py}, \texttt{cpsat\_check.py},
\texttt{cpsat\_decide.py}: validation.
\item \texttt{countbound.py}, \texttt{uc\_check.py}, \texttt{summary.py}: the bound $\CB$ and the
tables.
\item \texttt{reductions\_check.py}: \S\ref{sec:red}.
\item \texttt{mcut.c}, \texttt{evenfilter.c}, \texttt{oddfilter.c}, \texttt{random\_euler.py},
\texttt{multipartite.py}: generation.
\end{itemize}
\item Data: input graph6 files (the larger ones gzipped). Per-root tables \texttt{data/out/*.tab.gz} give
graph6, root, $\pr$, $n/2$ and status. For the two largest families (dense order $12$, $4$-regular
order $14$) the tables are per graph, \texttt{data/out/*.pg.gz}: graph6, the witness size at each root,
and a status letter per root. All witness dumps are regenerated by \texttt{run\_all.sh}, and
\texttt{data/SHA256SUMS} records the inputs.
\end{itemize}

\begin{thebibliography}{99}
\bibitem{AB} N.~Anto, M.~Basavaraju, Gallai's path decomposition for $2$-degenerate graphs,
\emph{Discrete Math.\ Theor.\ Comput.\ Sci.} 25:1 (2023); arXiv:2211.07159.
\bibitem{BBB} A.~Blanch\'e, M.~Bonamy, N.~Bonichon, Gallai's path decomposition in planar graphs,
arXiv:2110.08870v2 (2022); extended abstract in \emph{Extended Abstracts EuroComb 2021}, Trends in
Mathematics 14, 758--764.
\bibitem{BP} M.~Bonamy, T.~Perrett, Gallai's path decomposition conjecture for graphs of small maximum
degree, \emph{Discrete Math.} 342 (2019) 1293--1299; arXiv:1609.06257.
\bibitem{BS} F.~Botler, M.~Sambinelli, Towards Gallai's path decomposition conjecture, \emph{J.\ Graph
Theory} 97 (2021) 161--184; arXiv:1911.04546.
\bibitem{BSCL} F.~Botler, M.~Sambinelli, R.~S.~Coelho, O.~Lee, Gallai's path decomposition conjecture for
graphs with treewidth at most 3, \emph{J.\ Graph Theory} 93 (2020) 328--349; arXiv:1706.04334
(theorem numbers as in arXiv v1).
\bibitem{CW} Y.~Chu, Y.~Wang, Path decompositions of Eulerian graphs, \emph{Discrete Math.} 349 (2026)
114830; arXiv:2510.12806.
\bibitem{FHZ} G.~Fan, J.~Hou, C.~Zhou, Gallai's conjecture on path decompositions, \emph{J.\ Oper.\
Res.\ Soc.\ China} 11 (2023) 439--449.
\bibitem{Xie} M.~Xie, On Gallai's path decomposition conjecture, \emph{Filomat} 37:17 (2023) 5829--5834.
\bibitem{ZLH} J.~Zhang, Q.~Liu, Y.~Hong, Gallai's path decomposition conjecture for $K_5$-minor-free
graphs, arXiv:2611.21673 (2026).
\end{thebibliography}

\end{document}
